\begin{table}[htbp]
\centering
\scriptsize
\caption{Сравнение baseline-подходов и основных CatBoost-экспериментов на test-выборке}
\label{tab:practice-model-comparison}
\begin{tabularx}{0.98\textwidth}{L{0.24\textwidth} X r r r}
\toprule
Подход & Краткая характеристика & Accuracy & Macro F1 & Weighted F1 \\
\midrule
\texttt{most\_frequent} & Всегда выбирается самый частый класс обучающей выборки & 0.305068 & 0.051946 & 0.142623 \\
\texttt{last\_crop} & Следующая культура приравнивается к последней наблюдаемой культуре поля & 0.354321 & 0.172575 & 0.361135 \\
\texttt{Markov-1} & Одношаговая вероятностная модель переходов вида $P(\mathrm{target}\mid\mathrm{history}_3)$ & 0.557219 & 0.411371 & 0.539027 \\
Baseline CatBoost & История, регион, площадь и координаты поля & 0.699006 & 0.568485 & 0.689626 \\
CatBoost + extra features & Дополнительные признаки истории: разнообразие культур, индикаторы повторяемости, наличие отдельных групп культур & 0.699235 & 0.569473 & 0.689991 \\
CatBoost + transition graph & Гибридная оценка $\alpha p_{\mathrm{catboost}}+\beta p_{\mathrm{graph}}$ с лучшей конфигурацией $\alpha=0.9$, $\beta=0.1$ & 0.698714 & 0.564387 & 0.688729 \\
\bottomrule
\end{tabularx}
\end{table}
